\documentclass[12pt,a4paper]{article}

% ---------- Packages ----------
\usepackage[margin=1in]{geometry}
\usepackage{graphicx}
\usepackage{enumitem}
\usepackage{hyperref}
\usepackage{listings}
\usepackage{xcolor}
\usepackage{booktabs}
\usepackage{array}
\usepackage{titlesec}
\usepackage{fancyhdr}
\usepackage{tikz}

% ---------- Code listing style ----------
\lstset{
  language=Python,
  basicstyle=\ttfamily\small,
  keywordstyle=\color{blue}\bfseries,
  commentstyle=\color{green!60!black},
  stringstyle=\color{red},
  numbers=left,
  numberstyle=\tiny\color{gray},
  frame=single,
  breaklines=true,
  captionpos=b
}

% ---------- Header/Footer ----------
\pagestyle{fancy}
\fancyhf{}
\rhead{CSE 402 -- Operating System Lab}
\lhead{Lab Report 1}
\cfoot{\thepage}

% ---------- Document ----------
\begin{document}

% ===== TITLE PAGE =====
\begin{titlepage}
  \centering
  \vspace*{1.5cm}
  {\Large\bfseries CSE 402}\\[4pt]
  {\large University of Asia Pacific}\\[6pt]
  {\large Course Name: Operating System Lab}\\[12pt]
  {\LARGE\bfseries Lab Report 1}\\[8pt]
  {\Large Topic: FCFS Algorithm}\\[2cm]

  \begin{tabular}{@{}p{6cm}p{6cm}@{}}
    \textbf{Presented by,} & \textbf{Presented to,}\\[4pt]
    Md Mashrur Rahman Masfi & Atia Rahman Orthi\\
    Reg: 23101182 & Lecturer, UAP\\
    Section: D1 & \\
  \end{tabular}

  \vfill
  {\small GitHub: \url{https://github.com/MasfiRahman/CSE-402-Operating-System-Lab-}}
\end{titlepage}

% ===== PROBLEM STATEMENT =====
\section{Problem Statement}

In a multiprogramming environment the CPU must constantly decide which of the
waiting processes should execute next. This experiment implements the
\textbf{First Come First Serve (FCFS)} policy, where processes are dispatched
strictly in the order they enter the ready queue, and a running process keeps
the CPU until its burst finishes (non-preemptive scheduling).

The task is to write a Python program that accepts five processes, each
described by a \textit{Process ID}, \textit{Arrival Time} and
\textit{Burst Time}, and then:

\begin{itemize}[leftmargin=*]
  \item Arranges the processes in increasing order of arrival time.
  \item Simulates the non-preemptive execution on a single CPU to obtain each
        process's Completion Time (CT).
  \item Derives the Turnaround Time ($\text{TAT} = \text{CT} - \text{AT}$)
        and Waiting Time ($\text{WT} = \text{TAT} - \text{BT}$).
  \item Computes the average turnaround time and average waiting time.
  \item Prints the execution timeline (Gantt chart) so the result can be
        verified by hand.
\end{itemize}

% ===== SOURCE CODE =====
\section{Source Code}

% Replace the placeholder below with your actual Python code
\begin{lstlisting}[caption={FCFS Scheduling -- Python Implementation}]
# Paste your Python source code here
\end{lstlisting}

% ===== OUTPUT SCREENSHOT =====
\section{Output Screenshot}

% Uncomment and set the correct path for your screenshot
% \begin{figure}[htbp]
%   \centering
%   \includegraphics[width=0.9\textwidth]{output.png}
%   \caption{Console output of the FCFS program}
% \end{figure}

\textit{[Insert output screenshot here.]}

% ===== DISCUSSION =====
\section{Discussion}

The aim of this experiment was to implement and verify First-Come-First-Serve
(FCFS) scheduling on a small five-process workload. As expected from a
non-preemptive FIFO policy, the dispatch order followed arrival order, not
process numbering: P4 (AT\,=\,0) began immediately, and each subsequent CPU
assignment went to the process that had waited longest in the ready queue---P5,
then P2, P1, and finally P3.

The resulting Gantt timeline ($0 \to 3 \to 5 \to 6 \to 9 \to 11$) contains no
idle gaps, because at every completion instant at least one process was already
waiting. Consequently, the overall finish time of 11~units is exactly the sum
of the five bursts ($3+2+1+3+2=11$).

The computed metrics satisfy the standard identities
$\text{TAT} = \text{CT} - \text{AT}$ and $\text{WT} = \text{TAT} - \text{BT}$
for every process. Taking P3 as an example: it arrives at~5 and completes
at~11, giving a turnaround of~6; subtracting its 2-unit burst leaves a waiting
time of~4. Aggregating over all five processes yields an \textbf{average
turnaround time of 4.60} and an \textbf{average waiting time of 2.40}, in
exact agreement with the hand-calculated table, which validates the simulation.

At the same time, the trace clearly demonstrates FCFS's principal weakness: it
considers \emph{when} a process arrives but never \emph{how long} it will run.
P3, although one of the shortest jobs in the set, suffers the largest wait
(4~units) simply because it arrived last, and P1 spends 3~units queued behind
shorter jobs. Such convoy-like behaviour---short processes trapped behind
longer ones---inflates average waiting time and makes pure FCFS unsuitable for
interactive or response-sensitive systems. Even so, its simplicity,
predictability, and freedom from starvation explain why the FCFS principle
still appears as a building block inside more sophisticated policies such as
multilevel queues.

% ===== ADVANTAGES =====
\section{Advantages Observed}

\begin{itemize}[leftmargin=*]
  \item The entire scheduler reduces to a single FIFO queue, making
        implementation and testing straightforward.
  \item Execution order is fully predictable from arrival times, which
        simplifies debugging and manual verification.
  \item Service in arrival order guarantees every process eventually runs;
        starvation is impossible.
  \item Context switches occur only at termination or blocking, so scheduling
        overhead is minimal.
\end{itemize}

% ===== DISADVANTAGES =====
\section{Disadvantages Observed}

\begin{itemize}[leftmargin=*]
  \item Being non-preemptive, a single long CPU burst delays every process
        behind it.
  \item Average waiting time is rarely optimal and typically exceeds that of
        SJF, SRTF, or Round Robin.
  \item The convoy effect penalizes short jobs arriving behind long ones
        (seen here with P3).
  \item Response time is inadequate for time-sharing workloads.
  \item The policy offers no mechanism for priorities or deadlines.
\end{itemize}

% ===== CONCLUSION =====
\section{Conclusion}

The FCFS program arranged all processes in order of their arrival time and
simulated non-preemptive execution. The resulting completion, turnaround, and
waiting times matched both the manual calculations and the Gantt chart,
producing an average turnaround time of 4.60 and an average waiting time of
2.40. This exercise confirmed that FCFS is straightforward, deterministic, and
starvation-free. However, it also highlighted the algorithm's major drawback:
since burst time is completely disregarded, shorter processes can end up
waiting excessively behind longer ones---a phenomenon known as the
\textbf{convoy effect}. Because of this limitation, real-world operating
systems seldom use pure FCFS. Nevertheless, its core first-come-first-served
principle continues to serve as the building block for more advanced
scheduling policies, such as multilevel feedback queues.

\end{document}